\documentclass[../main.tex]{subfiles}

\begin{document}

\section{EDep-sim} \label{sec:edep-sim}
  EDep-sim (\textit{energy deposition simulation}) è un wrapper per
  GEANT4~\cite{Agostinelli_2003} un \textit{toolkit} per la simulazione del
  passaggio di particelle attraverso la materia e della deposizione di
  energia nel rivelatore.
  In input richiede file cinematici come i file di output di GENIE e file di
  geometria come quelli necessari a GENIE\@.
  Conoscendo la cinematica dell'evento e la geometria del rivelatore,
  EDep-sim propaga le particelle all'interno del materiale.\\
  Il risultato è restituito in output sotto forma di un file ROOT con la
  seguente struttura~\cite{edepsim_output}:
  \begin{itemize}
    \item File ROOT
    \begin{itemize}
      \item DetSimPassThru: Folder (TDirectoryFile) con informazioni su input
      \begin{itemize}
        \item gRooTracker: TTree di input formato
        rootracker~\cite[pag. 116]{Andreopoulos_2010}
        \item
        InputKinem: TTree con mappa tra eventi di input ed eventi simulati. Per
        ogni evento simulato riporta a quale file di input (inputFileNum) ed a
        quale evento (inputEntryNum) corrisponde
        \item
        InputFiles: Informazioni sui file di input (fileName, generatorName,
        treeName, filePOT, fileEntries)
      \end{itemize}
      \item EDepSimGeometry: TGeoManager con geometria
      \item EDepSimEvents: TTree con informazioni sulla simulazione degli eventi
    \end{itemize}
  \end{itemize}
  Gli oggetti di maggior interesse nell'ambito di questa tesi sono il
  gRooTracker che contiene l'output di GENIE e EDepSimEvents che invece
  contiene quello di EDep-sim.

\end{document}